% Transcription — fonds Alexandre Grothendieck, shelfmark 121, batch 2 (pages 21-40).
% Established from the facsimile archives/batches/121/batch-02.pdf (a copy of
% 121.pdf fetched through the project's relay,
% https://grothendieck-relay.onrender.com/source/121.pdf; PDF page k is page
% 20+k, no cover sheet), per the skill transcribe-grothendieck. The folder is
% seventy-two pages in four batches (1-20, 21-40, 41-60, 61-72), transcribed
% in parallel; this is the second.
%
% Pass: Opus 5.5 (claude-opus-5-5), 2026-09-26 — first pass, unchecked against the pages by a human.
%
% Numbering: \page{} follows the PDF count, which is the Montpellier footer
% (21-40). The archivists' pencil numbers bottom left run two ahead
% throughout this batch (23 on p. 21 ... 42 on p. 40), as batch 1 found.
% He paginates nothing on these leaves.
%
% Pages skipped: 24, 28, 37 and 39, blank cream separator sheets (only
% show-through on 37). Pages summarised: none; but the sketches of p. 40 are
% described in a \note rather than drawn, and p. 38, a sheet of formulas,
% has its layout described in notes.
% Pages transcribed: 21-23, 25-27, 29-36, 38, 40, all in his hand, blue ink.
%
% What the batch is about. Pp. 21-23 continue batch 1's set-theoretic
% approach: whether a set locally a union of leaves is globally one, the
% sets \overline{\mathcal{F}}_U attached to an open U, a counterexample on
% the real line (R^+ is not in \overline{\mathcal{F}} for the partition
% {0}, R^*), then a sheaf \underline{\mathcal{F}} of germs of subsets
% satisfying a)-d) and a local form of e), the stronger axiom e_n)
% (« \underline{F}-dim <= n »), and an NB (the sheaf need not be generated by
% its global sections: a 2-dimensional analytic torus without non-constant
% meromorphic functions); a condition f) is announced and left blank.
% Pp. 25-27: equisingularity of X along a non-singular Y (local product
% N x Y_U), eliminable (immobile) and nu-mobile singularities, the order of
% mobility nu(x), the open sets Mob^i(X) and the closed F_i(X), and
% Propositions A-D (the strata Y_i are non-singular of dimension i and X is
% equisingular along them; local product structure Y_i x (N, a) with a
% immobile; maximal germs of equisingular subvarieties; Y_nu(Y x N) =
% Y x Y_{nu-i}(N)). Pp. 29-36, written fair: « Équisingularité » in the
% piecewise-linear (or algebraic, analytic) context — equisingularity of a
% space and of a diagram of spaces along a subspace, equisingular
% filtrations, leaves (restricted and closed), incidence relation and
% boundary, equisingular / admissible / privileged partitions, the four
% canonical filtrations of a partition, Conjecture 1 (a coarsest privileged
% partition, « canonique ») and Conjecture 2 (for a finite diagram of proper
% maps), a Remark on proper maps sending leaves onto leaves as locally
% trivial fibrations, and a conjectural recursive construction of the
% privileged partition of a proper map f : X -> Y. Pp. 35-36 refer twice to
% « la précédente lettre à A'Campo » and end « ... proposée dans la 1re
% lettre : A'Campo » — these pages look like the draft of a letter of his to
% A'Campo, perhaps the 1974 letter of the inventory; they are his, carry
% mathematics, and are transcribed. P. 38: formulas for gluing sheaves on
% (Y <- Z -> X) — triples (F, G, alpha), i'_*, i'^*, p'_*, p'^*, q_!, and
% the strata X_{I,J} with the four operations and the dualising complex.
% P. 40 (a torn-out squared notebook sheet headed « Fujii Guruji
% Nichidatsu »): manifolds with boundary V_i^o indexed by a finite ordered
% set, boundary pieces ∂_j V_i^o fibred over V_j^o, and sketches.
%
% Notation fixed once for the batch, following batch 1: \mathcal{F},
% \overline{\mathcal{F}}_U, \mathcal{F}_X, conditions a)-f) as letters; his
% underlined \mathcal{F} of p. 22-23 (a sheaf) as \underline{\mathcal{F}};
% \mathrm{Mob}^i(X), F_i(X), Y_i, nu, \nu(x); \mathfrak{X} for his gothic X
% (a diagram of spaces); \Phi, \Phi_i for leaves and ordered sets of leaves;
% X | U for restriction, as written; \simeq for his vertical « 2 »-shaped
% isomorphism sign in the diagrams of pp. 25 and 29.
%
% Readings to check first: the interlinear additions of p. 21; the margin
% notes of pp. 22, 26, 33, 36; the struck lines and the end of p. 34; « ou
% immobile » p. 25; « divisible » p. 27; the indices on p. 40.

\documentclass[11pt,a4paper]{article}
% Le préambule commun à toutes les transcriptions du fonds Grothendieck.
%
% Il tient en une idée : **l'incertitude se compose, elle ne se résout pas.**
% Une transcription qui lisse un mot douteux a détruit la seule chose pour
% laquelle on la faisait — un lecteur doit pouvoir savoir, à chaque endroit,
% ce qui était sur le feuillet et ce qui a été ajouté. Les six macros qui
% suivent sont donc l'appareil critique, et non de la décoration.
%
% Les mêmes commandes sont reconnues par `scripts/render.mjs`, qui produit la
% vue de lecture du site. Une macro ajoutée ici doit l'être aussi là-bas,
% faute de quoi le rendu échouera bruyamment — ce qui est voulu : un rendu
% silencieusement faux est pire qu'un rendu absent.

\ProvidesPackage{grothendieck}[2026/08/08 Transcriptions du fonds Grothendieck]

\usepackage{amsmath,amssymb,amsthm}
\usepackage{mathrsfs}
\usepackage{stmaryrd}
% Les diagrammes commutatifs sont partout dans ce fonds : c'est la langue
% même de la géométrie algébrique de Grothendieck, et une transcription qui
% les rendrait en prose ne transcrirait rien.
\usepackage{tikz-cd}
% Français par défaut — c'est la langue des feuillets — mais l'anglais est
% chargé aussi : la traduction et le résumé sont des documents anglais, et la
% typographie française (espaces avant les deux-points) y serait une faute.
% Ces documents basculent par \selectlanguage{english} en tête de corps.
\usepackage[english,french]{babel}
\usepackage{fontspec}
\usepackage{geometry}
\geometry{a4paper,margin=25mm,marginparwidth=28mm}
\usepackage{marginnote}
\usepackage{xcolor}
% ulem, not soul. `soul`'s \st and \ul are built on a character-by-character
% scan that XeLaTeX's font machinery breaks: the document fails with an
% undefined control sequence rather than degrading. ulem does the same two jobs
% and is Unicode-engine safe. `normalem` stops it hijacking \emph, which the
% transcriptions use for Grothendieck's own emphasis.
\usepackage[normalem]{ulem}
\usepackage{hyperref}
\hypersetup{colorlinks=true,linkcolor=black,urlcolor=black,pdfborder={0 0 0}}

% --- Métadonnées du lot -----------------------------------------------------
% Reprises telles quelles par le script de rendu, qui les affiche en tête de
% la vue de lecture. Elles fixent ce que le fichier prétend couvrir : sans
% elles, une transcription flotte sans point d'attache dans un fonds de seize
% mille feuillets.

\newcommand{\folder}[1]{\gdef\grothFolder{#1}}
\newcommand{\batch}[1]{\gdef\grothBatch{#1}}
\newcommand{\leaves}[2]{\gdef\grothFirst{#1}\gdef\grothLast{#2}}
\newcommand{\foldertitle}[1]{\gdef\grothFoldertitle{#1}}
\gdef\grothFolder{?}\gdef\grothBatch{?}\gdef\grothFirst{?}\gdef\grothLast{?}\gdef\grothFoldertitle{}

% --- Le résumé --------------------------------------------------------------
%
% Il ouvre la lecture modernisée, sous le titre « Résumé », et il est composé
% en petit. Ce n'est pas un ornement : c'est le seul passage du document qui
% ne soit pas des mathématiques, et un lecteur doit pouvoir le reconnaître
% avant de l'avoir lu — puis le sauter s'il connaît déjà le sujet, ou s'y
% arrêter s'il ne le connaît pas. Le filet qui le suit marque la reprise du
% corps.

\newenvironment{resume}
  {\small\setlength{\parskip}{0.45em}}
  {\par\medskip\hrule\medskip}

% --- Mots-clés ---------------------------------------------------------------
%
% En anglais, en clôture du résumé de la lecture modernisée : le vocabulaire
% moderne sous lequel chercher ce que les feuillets construisent. Le manifeste
% du site (`npm run manifest`) les extrait pour étiqueter la cote entière —
% cette ligne est la source unique des tags, il n'y a pas de fichier à côté.
\newcommand{\keywords}[1]{\par\smallskip\noindent
  {\footnotesize\textsc{Keywords} --- \textit{#1}}\par}

% --- L'appareil critique ----------------------------------------------------

% Le numéro de feuillet, en marge. C'est l'ancre qui permet de régler un
% désaccord sur une formule en regardant *un* feuillet plutôt qu'une chemise
% de six cents. Le site s'en sert aussi pour tourner les pages du fac-similé
% quand on fait défiler la transcription.
\newcommand{\leaf}[1]{%
  \marginnote{\footnotesize\color{gray!70}\textsf{#1}}%
  \label{leaf:#1}\ignorespaces}

% Illisible : on ne devine pas. Un mot inventé qui se lit comme les autres est
% le pire résultat possible.
\newcommand{\ill}{\textcolor{red!60!black}{[\,\ldots\,]}}

% Lecture douteuse : le mot est donné, et signalé comme incertain.
\newcommand{\uncertain}[1]{\uline{#1}}

% Ajout de l'éditeur — une lettre manquante, un mot sous-entendu.
\newcommand{\add}[1]{\textcolor{blue!50!black}{[#1]}}

% Biffé par Grothendieck lui-même. Ce qu'il a rayé fait partie du document :
% une rature dit souvent où la pensée a changé de direction.
\newcommand{\struck}[1]{\sout{#1}}

% Note du transcripteur, sur l'état matériel du feuillet.
\newcommand{\note}[1]{\par\smallskip{\small\color{gray!60!black}\textit{[#1]}}\par\smallskip}

% Note marginale de Grothendieck, distincte du corps du texte.
\newcommand{\marginal}[1]{\par\smallskip\noindent\hspace{1em}%
  {\small\color{gray!60!black}#1}\par\smallskip}

% --- Titre ------------------------------------------------------------------

\AtBeginDocument{%
  \begin{center}
    {\large\bfseries Fonds Alexandre Grothendieck --- cote n\textsuperscript{o}~\grothFolder}\\[2pt]
    {\itshape\grothFoldertitle}\\[4pt]
    {\small feuillets \grothFirst--\grothLast\ (lot \grothBatch)}\\[2pt]
    {\footnotesize\color{gray!60!black}Transcription établie d'après le fac-similé numérisé
     par l'Université de Montpellier.\\
     Les lectures incertaines sont soulignées, les illisibles notées [\,\ldots\,],
     les ajouts entre crochets.}
  \end{center}
  \vspace{1em}}

\endinput


\folder{121}
\batch{2}
\pages{21}{40}
\dating{1974}
\watermark{Édition de démonstration}
\foldertitle{Topologie modérée : notes manuscrites (s.d.), lettre (1974)}

\begin{document}

\page{21}
On est ramené au cas où $U = X$, soit $A$ qui est loc.\ une réunion de
feuilles, l'est-il \emph{globalement} ? Il suffit de prouver que pour les
feuilles $X_i$, telles que $X_i \cap A \neq \emptyset$, on a $X_i \subset A$,
on \uncertain{a vu} que $X_i \cap A$ est une partie à la fois ouverte et
fermée de $A$\note{on attend plutôt « de $X_i$ » ; transcrit tel qu'écrit}.
C'est une condition locale sur $X$, or loc.\ sur $X$ $A$ est comme une
réunion de feuilles, OK !

Cela \emph{entraîne}, un peu plus gén., que pour \emph{tout} ouvert $U
\subset X$, si $A \in \overline{\mathcal{F}}_U$, alors les $A \cap X_i$
sont : à la fois ouverts et fermés dans $X_i \cap U$, donc des réunions de
comp.\ connexes de $X_i \cap U$. \add{Donc $A$ est réunion de comp.\
connexes \struck{\ill{}} des ensembles de la forme $X_i \cap
U$.}\note{addition interlinéaire, serrée entre deux lignes ; la barre
au-dessus de $\mathcal{F}_U$ porte, à sa première occurrence, une sorte de
flèche, et le signe est transcrit $\overline{\mathcal{F}}_U$ partout}
Inversement, supposons que $A$ soit une telle réunion, est-il vrai que $A
\in \overline{\mathcal{F}}_U$ ? \struck{Évidemment,} \struck{En particulier,
une comp.\ connexe}\note{« En particulier, une comp.\ connexe » est
récrit au-dessus de la ligne puis biffé avec elle}
\uncertain{une} composante connexe d'un $X_i \cap U$ est-il $\in
\overline{\mathcal{F}}_U$ ? \emph{Non}, exemple :

$X = U =$ droite numérique $\mathbb{R}$

\struck{$\mathcal{F} = \lbrace \emptyset, \mathbb{R}, \lbrace 0 \rbrace,
\mathbb{R}^{*} \rbrace$}

$X_0 = \lbrace 0 \rbrace$, $X_1 = \mathbb{R}^{*} = \mathbb{R} - \lbrace 0
\rbrace$\note{les trois lignes de l'exemple sont réunies par une accolade}


$\mathbb{R}^{+}$ \uncertain{n'est pas} $\in \overline{\mathcal{F}} =
\mathcal{F}_X$\note{les deux lettres sont écrites, la première avec une
barre, la seconde sans, au-dessus de l'indice $X$}.
Si on veut un exemple avec des $X_i$ connexes, $U \neq X$, on prend $X =
\widehat{\mathbb{R}} =$ \uncertain{$\mathbb{S}^1$}, \struck{\ill{}} $X_0 =
\lbrace 0 \rbrace$, $X_1 = X - \lbrace 0 \rbrace$, $U = \mathbb{R}$
\ldots\note{le signe après « $\widehat{\mathbb{R}} =$ » est surchargé ;
on lit une lettre ajourée d'exposant 1. Dans la marge gauche, à hauteur de
l'exemple, un signe entouré d'un ovale, illisible}

\page{22}
Nous étudions maintenant la condition f). \struck{Il est clair}
\add{On a vu} qu'elle est stable par restriction à un ouvert $\in
\mathcal{F}$, mais il est aussi clair qu'elle n'est pas stable par
restriction à un ouvert quelconque. Supposons, \uncertain{par} \uncertain{si}
$A \in \overline{\mathcal{F}}_U$ est loc.\ fermé dans $U$ (donc fermé dans un
ouvert plus petit $V \subset U$). Alors $A$ est loc.\ connexe, \struck{et
\ill{} !} donc ses comp.\ connexes sont ouvertes, si $A$ est loc.\ compact
(p.ex.\ si $X$ l'est) alors ses comp.\ connexes forment \uncertain{dans} une
famille loc.\ finie dans $A$ (donc dans $V$) \struck{\ill{}}. Elle
\add{\uncertain{est-elle}} \emph{loc.\ finie dans $U$} ?
\note{« Il est clair » est souligné puis biffé, « On a vu » écrit
au-dessus ; au-dessus de « Elle », deux signes (\uncertain{« est »} et un
signe d'égalité triple) ; la fin de la phrase est soulignée}

\marginal{\ill{} \uncertain{la top.\ de} \ill{} \uncertain{engendrée}
\ill{} \uncertain{ouverts} \ill{}}\note{note en biais dans la marge
gauche, à hauteur des lignes 8 à 10, presque entièrement illisible}

\struck{Ceci n'\uncertain{aurait} \ill{} \ill{} \ill{} \ill{} $U$, \ill{}
\ill{} $A = \ill{}$, avec $Z \in \mathcal{F}$, $Z$ loc.\ fermé,}
\ill{} $\forall x \in U$, $\exists$ voisinage ouvert $V \in
\mathcal{F}_U$, tel que $A \cap V \in \mathcal{F}$, et $A \cap V$ réunion
finie de comp.\ connexes.\note{les trois premières lignes de ce paragraphe
sont barrées de trois longs traits obliques, qui ne touchent pas la fin ;
un double trait horizontal le sépare de ce qui suit}


Soit $\underline{\mathcal{F}}$ un sous-faisceau du faisceau des germes de
parties de $X$. On suppose $\underline{\mathcal{F}}$ satisfait a) b) c) d),
et \struck{\ill{}} \add{de plus}
\begin{itemize}
  \item[e)] Soit $U$ un ouvert de $X$, et $X_1 \supset X_2
  \supset \cdots \supset X_n \supset \cdots$ des fermés de $U$, $\in
  \Gamma(U, \underline{\mathcal{F}})$, tels que $X_{n+1}$ rare dans $X_n$.
  Alors la suite des $X_n$ est loc.\ finie sur $U$, donc $\bigcap X_n =
  \emptyset$
\end{itemize}
\note{le premier mot de la condition, souligné, est biffé ; on y lit
« f\ldots\ », et « e) » est écrit au-dessus. Les deux $\mathcal{F}$
soulignés de cette page (faisceau) sont transcrits
$\underline{\mathcal{F}}$}

\page{23}
(dans \uncertain{pour} toute partie quasi-compacte $K$ de $U$, $\exists\, i =
i(K) \in \mathbb{N}$ tel que $X_i \cap K = \emptyset$). Alors $\mathcal{F} =
\Gamma(X, \underline{\mathcal{F}})$ satisfait à a) b) c) d) \struck{e)}, et
aussi à e) si $X$ \uncertain{quasi-compact}. Plus généralement, si $A \in
\overline{\mathcal{F}}$, alors $\mathcal{F}_A$ satisfait à a) b) c) d)
\struck{e)} (et satisfait à e) si ???).
\add{\uncertain{$=$} ensemble des parties de $A$ qui sont loc.\ \ill{} par
des sections de $\underline{\mathcal{F}}$}\note{addition interlinéaire
entourée, au-dessus de « $\mathcal{F}_A$ satisfait », reliée à
$\mathcal{F}_A$ ; lecture incertaine}
Ici il est plus commode de partir de l'axiome suivant plus fort que e) :
\begin{itemize}
  \item[e$_n$)] $\exists$ un entier $n$ tel que si $X_0 \supset X_1 \supset
  \cdots \supset X_n \supset X_{n+1}$ \add{\uncertain{dans} $U$}, comme dans
  e) ci-dessus, alors on a $X_{n+1} = \emptyset$. (On dit que $X$ est de
  $\underline{\mathcal{F}}$-dim $\leq n$)
\end{itemize}
\note{un trait vertical dans la marge gauche marque l'énoncé}
Alors tous les $\mathcal{F}_A$ ($A$ \struck{ouvert} fermé dans $X$),
satisfont à e$_n$), donc e).

\emph{NB} $\underline{\mathcal{F}}$ est $\supset$ \uncertain{pour} le
faisceau \emph{engendré} par $\mathcal{F} = \Gamma(X,
\underline{\mathcal{F}})$ et l'inclusion peut être stricte. \emph{Ex} :
germes de parties analytiques d'un tore analytique de dim 2 sans fonctions
méromorphes non constantes\ldots\note{« $\supset$ » : le signe après
« est » est petit et peu net}

Nous \struck{supposons} \add{envisageons} maintenant une condition :

f)
\note{la condition f) est annoncée mais laissée en blanc ; la page s'arrête
là}

\page{25}
1.\ $X$ espace top.

$Y$ sous-espace « non singulier » resp.\ variété

$X$ est dit « \emph{équisingulier le long de $Y$} » si $\forall y \in Y$,
\struck{\ill{}} $\exists$ espace pointé $(N, a)$, \struck{et} un voisinage
ouvert $U$ de $y$ dans $X$, et un homéomorphisme

\begin{tikzcd}
U \arrow[r, "\simeq"] & N \times Y_U \\
Y_U \arrow[u, hook] \arrow[r, "\simeq"] & a \times Y_U \arrow[u, hook]
\end{tikzcd}

($Y_U = Y \cap U$), avec commutativité.\note{le diagramme est redessiné :
sur la page, « avec commutativité : » précède la ligne du bas, et les deux
flèches verticales sont tracées vers le haut, un crochet à leur pied ;
« can » est écrit sous la flèche du bas. Dans $N \times Y_U$, un premier
$X$ est biffé}

On dit que $x \in X$ est une \emph{singularité éliminable}
(\struck{ou \ill{}} \add{ou immobile}) de $X$ si $\forall Y \subset X$, $Y$
sous-espace \struck{\ill{}} non singulier de $X$, $Y \ni x$, $X$
équisingulier \struck{\ill{}} le long de $Y$, on a $\dim_x Y =
0$.\note{« ou immobile » : le mot entre parenthèses, souligné, est biffé
et récrit au-dessus ; lecture incertaine du mot biffé}

On dit que $x \in X$ est une \emph{singularité $\nu$-mobile}
\struck{\ill{}} si $\exists\, Y \subset X$ sous-espace, $Y$ non sing.,
$Y \ni x$, $X$ équisingulier le long de $Y$, $\dim_x Y \geq \nu$ (N.B.\ on
peut alors choisir $Y$ avec $\dim Y = \nu$).

\struck{Ainsi, les singularités éliminables}

On dit que $x$ a un \emph{ordre de mobilité $\nu$} si $x$ est $\nu$-mobile
mais non $(\nu + 1)$-mobile. (Ainsi les singularités éliminables
\add{ou immobiles} sont celles qui sont d'ordre de mobilité $0$).
\struck{Donc les points} \add{\struck{Les points qui \ill{} un ordre de
mobilité}} \struck{l'un des points \ill{} ordre de \ill{} mobiles}
$\nu \geq \nu_0$ ($\nu_0$ donné) sont les pts $\nu_0$-mobiles qui ont un
ordre de mobilité. (Mais si $X$ \struck{est} loc.\ de dim finie, tout pt
$x$ de $X$ a nécessairement un ordre de mobilité $\nu(x)$, qui est le sup
des dim en $x$ des ss-espaces de $X$ \struck{\ill{}} non singuliers
\struck{qui sont le} le long desquels $X$ est
équisingulier)\note{interligne très chargé ; le début de la phrase
« \ldots\ sont les pts $\nu_0$-mobiles » est reconstitué à travers deux
essais biffés}

\page{26}
\struck{Supposons que $\nu(x)$ existe pour $\forall x \in X$ (p.ex.\ $X$
loc.\ de dim finie)}

L'ens.\ des pts $x \in X$ tels que \struck{$\nu(x) \geq i$} \add{$x$ soit
$i$-mobile dans $X$} \add{est un} \emph{ouvert} $\mathrm{Mob}^i(X)$ --- on
trouve \uncertain{ainsi} un \add{système} décroissant
\[
  \mathrm{Mob}^0(X) = X \supset \mathrm{Mob}^1(X) \supset \mathrm{Mob}^2(X)
  \supset \cdots
\]
\add{dont l'intersection est vide si $\forall x \in X$, $\nu(x)$ défini
(p.ex.\ $X$ loc.\ de dim finie)}. Les complémentaires forment des
\struck{\ill{}} ensembles
\[
  F_0(X) \subset F_1(X) \subset F_3(X) \subset \cdots \qquad
  (F_i(X) = X - \mathrm{Mob}^{i+1}(X))
\]
\marginal{réunion des $F_i(X)$ est $X$ ssi $\forall x \in X$, $\nu(x)$
défini) p.ex.\ $X$ loc.\ de dim finie}\note{écrit à droite de la suite des
$F_i$ ; l'indice $F_3$ est sic, pour $F_2$}

Pour tout $X$ fixé, on se pose la question \struck{suivante :} de la
validité des propositions « suivantes » (\uncertain{où} $F_{-1}(X) =
\emptyset$)

\emph{Prop A} $\forall i \geq 0$, \add{$Y_i =$} \struck{\ill{}} $F_i(X) -
F_{i-1}(X)$ \struck{$= \nu^{-1}(\lbrace i \rbrace)$} (ens.\ des pts d'ordre
de mobilité $i$) est un sous-espace non \struck{singulier} partout de
dimension $i$, et $X$ est équisingulier le long de $Y_i$.
\note{au-dessus de la ligne, « $Y_i =$ » est ajouté ; le signe final
« $= \nu^{-1}(\lbrace i \rbrace)$ » est surchargé et paraît biffé ; « non
singulier » : le second mot est souligné d'un trait qui le barre à demi}

\marginal{Appelons « \uncertain{vraies} \uncertain{variétés} » les espaces
loc.\ de dim finie qui \uncertain{vérifient} \ill{}}\note{note verticale
dans la marge gauche, reliée par une accolade à la Prop A ; lecture très
incertaine}

\emph{Prop B} Localement au voisinage de chaque pt de $Y_i$, on a $(X,
Y_i) \simeq Y_i \times (N_\xi, a)$, où \ill{} est un espace topologique dont
$a$ est une singularité \emph{immobile}. [Le type d'homéomorphie local de
$(N, a)$ ne dépend que de la composante connexe du point envisagé de
$Y_i$]\note{Dans la marge gauche, reliée par
une accolade à la Prop B : « vrai \uncertain{pour les} \ill{} + \ill{}
$\forall$ \ill{} $N$ \struck{\ill{}} \uncertain{divisible} ? »}

\emph{Prop C} Pour $\forall$ \struck{$x$} $\in Y_i$, $(Y_i, y) \mapsto$ le
plus grand germe de sous-variété $\ni y$ de $X$ en $y$ le long de laquelle
$X$ soit équisingulier. Tous les « \uncertain{autres} » $Y$ correspondent
$\xrightarrow{1-1}$ aux germes \add{en $y$} de sous-variétés régulières
\uncertain{plongées} de \add{$Y_i$} \struck{\ill{}}.
\note{« $\xrightarrow{1-1}$ » : le 1-1 est écrit au-dessus d'une flèche
en crochet ; la fin de la ligne, sous « de », est biffée}

\page{27}
\emph{Prop D} Soit $Y$ \struck{\ill{}} une variété topologique partout de
dim $i$, $N$ un espace topologique. Alors \struck{\ill{}} $\forall (y, n)
\in Y \times N$, \struck{$(y, n)$ est \ill{}} \struck{$(y, n)$ est
$i$-mobile.} si $\nu \geq i$, $(y, n)$ \add{$\in Y \times N$} est
$\nu$-mobile ssi $n \in N$ est $(\nu - i)$-mobile. Donc
\[
  Y_\nu(Y \times N) = Y \times Y_{\nu - i}(N).
\]
\marginal{$N$ \uncertain{divisible} $\Longleftrightarrow$ $Y \times N$
\uncertain{divisible} ($Y \neq \emptyset$)}
\marginal{Prop B $\Leftarrow$ (\uncertain{sous réserve} de Prop A)}
\note{deux notes dans la marge gauche, séparées de l'énoncé par un trait
vertical ; « divisible » est d'une lecture douteuse aux deux places}

\page{29}
\section*{Équisingularité}
\note{titre souligné, de sa main}

\struck{\ill{}} Le contexte des « espaces » est celui des espaces
topologiques \add{(ou algébriques, ou analytiques)} linéaires \add{par
morceaux}, \struck{\ill{}} \struck{espaces algébriques, par} Le mot
« espace » et « variété » désignent les \struck{\ill{}} objets et
morphismes du contexte choisi.\note{« (ou algébriques, ou analytiques) »
est écrit au-dessus de la ligne et souligné ; « par morceaux » est ajouté
dans la ligne suivante, au-dessus d'un trait}

Soient $X$ un espace, $Y$ un sous-espace loc.\ fermé. $X$ est dit
\emph{équisingulier le long de $Y$} si $\forall y \in Y$, $\exists$ ouvert
$X'$ de $X$, et espace pointé $(Z, a)$, tels que l'on ait un diagramme
commutatif

\begin{tikzcd}
X' \cap Y = Y' \arrow[r, hook] \arrow[d, "\simeq"] & X' \arrow[d, "\simeq"] \\
Y' \times \lbrace a \rbrace \arrow[r, hook] & Y' \times Z
\end{tikzcd}

\note{les deux flèches verticales portent un signe en forme de 2, qu'on
lit comme un $\simeq$ couché ; au lieu de $X'$ il avait d'abord écrit une
autre lettre, surchargée}

Plus généralement, soit $\mathfrak{X} = (X_i)_{i \in I}$ un diagramme
d'espaces ($I$ étant un type de diagrammes, p.ex.\ défini par un ens.\
ordonné --- c'est le seul cas que nous envisagerons), soit $i_0 \in I$ et
$Y$ un sous-espace loc.\ fermé de $X_{i_0}$. On dit que $\mathfrak{X}$ est
\emph{équisingulier le long de $Y$} si, désignant par $J$ le sous-diagramme
de $I$ formé de tous les $j \in I$ pour lesquels \add{ou bien $j = i_0$ ou
bien} $\exists$ flèche $j \to i_0$, $\forall y \in Y$ $\exists$ voisinage
ouvert $X'_{i_0}$ de $y$ dans $X_{i_0}$, \add{d'espaces} un diagramme
$(Z_j)_{j \in J}$, avec $Z_{i_0}$ muni d'un point $a \in Z_{i_0}$, et un
homomorphisme de diagrammes ($Y' = Y \cap X'_{i_0}$)
\[
  (Z_j \times Y')_{j \in J} \longrightarrow (X_j)_{j \in J}
\]
tels que pour tout $j \in J$, $Z'_j \times Y' \to X_j$ soit un plongement
ouvert, et que $\exists$ diagramme comm.

\begin{tikzcd}
Y' \arrow[r, hook] \arrow[d, "\simeq"] & X'_{i_0} \arrow[d, "\simeq"] \\
\lbrace a \rbrace \times Y' \arrow[r, hook] & Z_{i_0} \times Y'
\end{tikzcd}

\note{« d'espaces » est écrit au-dessus de « diagramme », avant
$(Z_j)$ ; $Z_{i_0}$ est écrit une fois $Z'_{i_0}$, et $Z'_j$ dans
« $Z'_j \times Y' \to X_j$ » : transcrit tel qu'écrit. Le « et que »
ajouté en fin de phrase est de lecture incertaine}

\page{30}
Soit $X$ un espace, $(X_i)_{i \geq 0}$ une filtration croissante (pour fixer
les idées) de $X$ par des sous-espaces fermés. On dit que la
\emph{filtration est équisingulière} si $\forall i \geq 0$, le diagramme
formé de $(X_0 \to X_1 \to \cdots \to X)$ est équisingulier le long de
$X_i - X_{i-1} = V_i$ (NB on pose $X_{-1} = \emptyset$).

Appelons « \emph{feuilles} (restreintes) » de la filtration les
composantes connexes d'un $V_i = X_i - X_{i-1}$. Les feuilles restreintes
forment donc une \emph{partition} de $X$ si $X$ est réunion des $X_i$, ce
qu'on supposera désormais.\note{un crochet dans la marge gauche marque le
début de ce paragraphe}

\emph{Proposition} Si la \struck{condition} filtration est
équisingulière, alors l'adhérence de toute feuille est une réunion de
feuilles.

(Si $X$ est compact, et les $X_i$ loc.\ connexes --- ce qui est vrai dans
le contexte \uncertain{an.\ ou alg.}\ ou aussi « par morceaux » --- alors
il s'agit bien sûr d'une réunion \emph{finie} de feuilles\ldots)

Les adhérences des feuilles seront appelées \emph{feuilles fermées}. Il y a
corr.\ biunivoque via $F \mapsto \overline{F}$ entre feuilles (restreintes
--- on dira aussi \uncertain{les} appelées « ouvertes ») et feuilles
fermées. Les parties \emph{fermées} de $X$ qui sont réunions de feuilles
sont aussi celles qui sont réunions de feuilles fermées. Toute intersection
et toute réunion finie de telles parties est une

\page{31}
partie du même type. Par exemple toute intersection de feuilles fermées
est une réunion de feuilles fermées.

La relation d'inclusion entre feuilles fermées induit une relation d'ordre
sur l'ens.\ des feuilles, appelée « relation d'incidence ». La réunion des
feuilles incidentes à une feuille donnée $F$ et distinctes d'icelle n'est
autre que $\overline{F} - F$ ; on l'appelle le \emph{bord} de la feuille
$F$ ; c'est donc une partie fermée, réunion de feuilles fermées.

\emph{Proposition} Pour que la filtration donnée de $X$ soit
équisingulière, il f.\ et s.\ qu'elle satisfasse les 2 conditions :
\begin{itemize}
  \item[a)] \struck{$\forall$} feuille $F$, $\overline{F}$ est une réunion
  de feuilles
  \item[b)] Le diagramme formé de $X$ et de toutes les feuilles fermées
  (avec les \struck{applications} \add{\uncertain{inclusions}}
  d'inclusion) est \struck{\ill{}} \add{équisingulier} le long de chaque
  feuille.
\end{itemize}
\note{au-dessus de « applications » biffé, un mot surchargé, peut-être
« inclusions » ; « équisingulier » est écrit au-dessus d'un mot biffé}

On dit qu'une partition $(F_i)_{i \in I}$ \struck{d'un} \add{en
« feuilles » $F_i$} \add{d'un} espace $X$ \struck{est équisingulière} est
une \emph{partition équisingulière} si 1°) chaque $F_i$ est la différence de
deux sous-espaces fermés (i.e.\ $\overline{F_i}$ et $\overline{F_i} - F_i =
bF_i$ sont des sous-espaces fermés) 2°) Chaque $\overline{F_i}$ est une
réunion de feuilles et 3°) le diagramme formé par \add{$X$,} les
$\overline{F_i}$ et leurs inclusions est équisingulier le long de chaque
feuille.

\marginal{Si 1° et 2° sont satisfaits, on dira que la partition est
\uncertain{admissible}}\note{dans la marge gauche, en biais, reliée par
une accolade aux conditions 1°) et 2°)}

\page{32}
\emph{NB} Deux filtrations \struck{\ill{}} \add{équisingulières} de $X$
peuvent très bien définir des \struck{\ill{}} partitions identiques de $X$
en feuilles. C'est la partition en feuilles qui est la plus importante,
non la filtration. P.ex.\ dans le contexte « par morceaux », où on dispose
donc d'une notion de « dimension » et de « codimension », il y a au moins 4
façons canoniques, \struck{\ill{}} \add{distinctes} « en général », pour
décrire une partition équisingulière donnée à l'aide de filtrations : 1°)
$X_i$ = réunion des feuilles de dimension $\leq i$, 2°) $X_0$ = réunion des
feuilles minimales \ldots\ $X_i$ = réunion des feuilles minimales parmi
celles qui sont $\not\subset X_{i-1}$, 3° et 4° définitions « duales » de
1° et 2°, donnant des filtrations \emph{décroissantes}.

Pour la suite, on préfèrera utiliser le langage des espaces à partitions
équisingulières, plutôt que des espaces filtrés équisinguliers.

\struck{\emph{Prop} Si la partition donnée de $X$ est équisingulière, \ill{}
\ill{} les \ill{} fermés de $X$ qui sont réunions de feuilles, \ill{} la
partition \ill{} équisingulière.}\note{deux lignes serrées, encadrées d'un
crochet et biffées d'un long trait}

Soit $(X_i)_{i \in I}$ un diagramme d'espaces, et supposons que tout $X_i$
est muni d'une partition \struck{admissible}. On \struck{dit que la
définit la} \add{dit que} \struck{système de partitions équisingulières}
\add{le système de} partitions est \emph{admissible} si
\begin{itemize}
  \item[a)] $\forall i \in I$, la partition donnée de $X_i$ est admissible
  \item[b)] $\forall$ flèche $i \to j$, l'application $X_i \to X_j$ applique
  une feuille de $X_i$ dans une feuille (évidemment unique déterminée) de
  $X_j$.
\end{itemize}

\page{33}
Lorsque cette condition est satisfaite, on définit un nouveau diagramme
(dont le type se déduit de façon évidente du type de diagramme $I$), des
ensembles ordonnés par incidence \struck{des} $\Phi_i$ des feuilles des
espaces $X_i$, et des morphismes d'ensembles ordonnés \struck{\ill{}}
$\Phi_i \to \Phi_j$ définis par les \struck{flèches} applications $X_i \to
X_j$ définies par les flèches du diagramme $I$ --- d'où un diagramme
\add{$(\Phi_i)_{i \in I}$} \struck{de} type $I$ d'ensembles ordonnés\ldots),
par les applications $F_{i,\alpha} \to F_{j,\beta}$ entre feuilles
\emph{fermées} induites par les $X_i \to X_j$. Ceci dit, le
\struck{système} \add{système} \struck{diagramme} des partitions des $X_i$
est dit \emph{équisingulier} (relativement au diagramme donné $\mathfrak{X}
= (X_i)_{i \in I}$) \struck{s'il est admissible, et si} \add{\uncertain{
prendre}} si le diagramme \add{des} \struck{\ill{}} applications entre
feuilles fermées est équisingulier le long de toute feuille (restreinte) de
tout $X_i$.\note{« \uncertain{prendre} » : mot écrit au-dessus de
« si » ; peut-être « prenant »}

\marginal{Un système de partitions est dit « privilégié » si \ill{}
\uncertain{équisingulier} \ill{} les feuilles \ill{} des variétés (sans
bord)}\note{note verticale dans la marge gauche, reliée par une accolade à
la définition ; lecture fragmentaire}

On suppose maintenant que le contexte soit linéaire (ou algébrique, ou
analytique) \emph{par morceaux}.

\emph{Conjecture 1} Tout espace $X$ admet une partition
\struck{équisingulière} \add{privilégiée}. Parmi toutes les partitions
\struck{équisingulières} \add{privilégiées} de $X$, il en est une plus
grossière que toutes les autres (appelée « canonique »).

\page{34}
Cela ne suffit pas pour bien des questions : on aura besoin de partitions
privilégiées qui \struck{\ill{}} soient \add{admissibles} \add{\uncertain{(pas
grossières)}} plus fines qu'une partition donnée, ou telles que
\struck{données} \add{les $X_i$} par des sous-espaces fermés donnés de $X$
en nb fini, chacun soit réunion de feuilles et que la filtration induite
soit \struck{équisingulière} \add{privilégiée}. Ou pour une application
donnée $X \to Y$, on voudrait trouver simultanément des partitions
privilégiées compatibles avec cette application. Tous ces cas semblent
couverts par l'énoncé conjectural suivant :\note{l'interligne de la
deuxième ligne est chargé : « admissibles » et « (pas grossières) » sont
écrits au-dessus d'un mot biffé ; lecture incertaine}

\emph{Conjecture 2} Soient $I$ un ens.\ ordonné \emph{fini},
$(X_i)_{i \in I}$ un diagramme de type $I$ d'espaces, avec les applications
$X_i \to X_j$ \emph{propres}. Alors il existe des partitions privilégiées
de ce diagramme, et il y a parmi elles une moins fine (appelée comme de
juste « canonique »).

\emph{Remarque} Soit une \struck{\ill{}} partition admissible pour un
diagramme $(X_i)_{i \in I}$. Si $X_i \xrightarrow{f} X_j$ dans ce diagramme
est \emph{propre}, on voit aussitôt que pour \struck{\ill{}} \add{une}
feuille $F_\alpha$ de $X_i$, s'appliquant \add{donc} dans une feuille
$\Phi$ de $X_j$, l'image est $\Phi$ tout entier (c'est un\ldots\
\add{ou l'image de $\ldots$ dans $\Phi$}) \struck{\ill{}} seulement un
ouvert de $\Phi$) \struck{\ill{}} $F \cap f^{-1}(\Phi)$ est à la fois
ouvert \add{\uncertain{à} \uncertain{plein}}\note{fin de page très
raturée : les parenthèses et les additions interlinéaires se chevauchent,
et le début de la dernière ligne est biffé d'une rature serrée ; la phrase
se poursuit p.~35}

\page{35}
équisingularité) et fermé (par propreté) donc égal à $\Phi$ qui est connexe,
\struck{et d'autre part} \add{donc} \add{($\subset \overline{F} \cap
f^{-1}(\Phi)$)} $\Phi - f(F)$ est une réunion d'images de feuilles
\struck{\ill{}} de $X_i$, donc ouvert dans $\Phi$, donc $f(F)$ est à la
fois ouvert et fermé dans $\Phi$, donc $= \Phi$ ] De plus, \struck{bien}
\struck{\ill{}} comme $F$ est fibré loc.\ trivial sur l'ouvert $f(F)$ de
$\Phi$, \struck{et que} (équisingularité) il est en fait fibré loc.\ trivial
sur $\Phi$ tout entier.\note{« égal » est souligné ; le crochet fermant
répond à un crochet ouvert qui n'est pas sur la page précédente}

Mieux, $f^{-1}(\Phi)$ \struck{est fibré avec \ill{}} \add{\ill{}}
\emph{partition induite} est un fibré loc.\ trivial sur $\Phi$. Comme $F$
\struck{est} et $\Phi$ sont des variétés, il s'ensuit que les fibres de $F$
sur $\Phi$ doivent aussi être des variétés\ldots

Dans la \uncertain{précédente} lettre à A'Campo, j'ai donné une
description conjecturale de la « partition privilégiée canonique » d'un
espace $X$, en termes de points « immobiles ». Voici une description
conjecturale de la « partition privilégiée » d'une application propre
\struck{\ill{}}
\[
  f : X \to Y
\]
Pour simplifier, on suppose $Y$ compact, donc $X$ compact. On
\struck{procède par récurrence sur $\dim Y$ pour \ill{} filtration
décroissante de} \add{\uncertain{procède par}} la construction. Considérons
le plus grand ouvert $U$ de $Y$ tel que 1°) $U$ est non singulier, 2°) $X
| U \to U$ est un fibré loc.\ trivial. Il faut prouver que $Y - U = Y'$ est
un sous-espace (de la catégorie « par morceaux » envisagée) \emph{rare}
(donc $\dim Y' < \dim Y$). Les composantes connexes de $Y - Y'$ seront les
feuilles\note{le nom est écrit « A-Campo ». Les deux lignes biffées après
« On » ne se lisent qu'en partie}

\page{36}
ouvertes de la \struck{\ill{}} \add{partition} cherchée de $Y$. Les feuilles
de $X$ au-dessus de celles-ci sont simplement les feuilles de la partition
\add{privilégiée} canonique (intrinsèque) de $V = X | U$ (loc.\ sur $U$,
c'est le produit de $U$ avec une feuille canonique d'une fibre). Si $F_i$
sont ces feuilles, il faudra bien sûr que les $\Phi'_i = \overline{F_i} - V
\cap \overline{F_i} = \overline{F_i} | Y'$ soient des réunions de feuilles
de $X$. Mais si $X' = X | Y' = X - V$ on est ramené à trouver la
\struck{\ill{}} \add{partition canonique} pour $X' \to Y'$, mais $X'$ étant
muni de plus des sous-espaces $\Phi'_i$ --- et de leurs intersections
finies, \uncertain{sans quoi faire}. On gagne par hyp.\ de récurrence, à cela près
qu'il aurait fallu partir non avec $X \to Y$, mais avec en plus une famille
finie de sous-espaces fermés donnés de $X$. Pour paraphraser dans ce cas
plus général la construction précédente, il faut simplement (remplaçant $X$
par $V$) pour un espace $X$ muni d'une famille \add{finie} \struck{\ill{}}
de sous-espaces fermés, définir la \struck{\ill{}} partition privilégiée
canonique (i.e.\ la plus grossière) pour ce système. Or la construction
peut probablement se faire en paraphrasant celle proposée dans la
\add{1re} lettre : A'Campo.

\marginal{Il faudra \uncertain{prendre} pour les $\overline{F_i}$
(\uncertain{donc} $\overline{F_i}$) \ill{} \ill{} des « variétés » \ill{}
\uncertain{catégorie}}\note{note verticale dans la marge gauche, en trois
colonnes, rattachée par un signe à « feuilles de $X$. » ; lecture
fragmentaire. Le mot abrégé avant « : A'Campo » se lit « lettre », surmonté
de « 1re »}

\page{38}
\note{feuillet de formules sans texte suivi, en deux colonnes ; on les donne
dans l'ordre de la page, colonne de gauche d'abord}
\[
  \mathfrak{X} = \bigl( Y \xleftarrow{\;p\;} Z \overset{i}{\hookrightarrow} X
  \bigr), \qquad p^{*}(\mathcal{F}) \xrightarrow{\;\alpha\;}
  i^{*}(\mathcal{G}), \qquad (\mathcal{F}, \mathcal{G}, \alpha)
\]
\note{$\mathcal{F}$ est écrit au-dessus de $Y$, $\mathcal{G}$ au-dessus de
$X$ ; la flèche $p$ est barrée d'un petit trait vertical, et un point
marque la flèche $i$}

\struck{\ill{}}
\begin{align*}
  i'_{*}(\mathcal{F}) &= (\mathcal{F}, 0, 0) \\
  i'^{*}(\mathcal{F}, \mathcal{G}, \alpha) &= \mathcal{F} \\
  p'_{*}(\mathcal{G}) &= (p_{*} i^{*} \mathcal{G}, \mathcal{G},
  \mathrm{can} : p^{*} p_{*}(i^{*}\mathcal{G}) \to i^{*}(\mathcal{G})) \\
  p'^{*}(\mathcal{F}, \mathcal{G}, \alpha) &= \mathcal{G} \\
  q_{!}(\mathcal{G}) &= (0, \mathcal{G}, 0) \\
  q^{!}(\mathcal{G}) &=
\end{align*}
\note{dans $i'^{*}$, l'argument est surchargé ; dans $q_!$, le $0$ est
écrit sur une autre lettre ; $q^!$ reste inachevé}

\begin{tikzcd}[column sep=small, row sep=small]
 & Z \arrow[dl, "p"'] \arrow[dr, dashed, hook, "i"] & & \\
Y \arrow[dr, "i'"'] & & X \arrow[dl, "p'"] & X - Z \arrow[l, hook'] \arrow[dll, bend left=30, "q"] \\
 & \mathfrak{X} & &
\end{tikzcd}

\note{schéma en haut à droite de la page, redessiné ; $i'$ est une flèche
en crochet sur la page}

\begin{align*}
  & X_{I,J} = X_I - X_J, \qquad \Sigma \supset I \supset J \\
  & X_{I,J} \xrightarrow{\;\alpha_{IJ/I'J'}\;} X_{I',J'} \\
  & \alpha_{*},\ \alpha^{*},\ \alpha_{!},\ \alpha^{!} ; \qquad
  K_{X,\Lambda} \ \text{complexe dualisant}
\end{align*}
\note{sous « $I \supset J$ », une ligne peu lisible : « $J$
\uncertain{une antichaîne} \uncertain{pour} \ill{} \ill{} » ; les deux
dernières lignes sont réunies par une accolade}

\[
  \partial X_\alpha \xrightarrow{\;p_\alpha\;} \cdots, \qquad
  \partial X_\alpha \xrightarrow{\;i_\alpha\;} X_\alpha
\]
\note{au milieu de la page, $i_\alpha$ est une flèche verticale de
$\partial X_\alpha$ vers $X_\alpha$ ; le but de $p_\alpha$ est biffé et
illisible, \struck{\ill{}}}

\page{40}
\note{feuillet de cahier à petits carreaux, détaché d'une spirale ; en tête,
de sa main : « Fujii Guruji Nichidatsu ». La page est un brouillon de
formules et de croquis ; on transcrit les formules, et les croquis sont
décrits en fin de page}

$I$ ens.\ ordonné fini

$\forall i \in I$ \quad $V_i^\circ$ variété à bord [bord vide si $i$
minimal)

$\forall (i,j) \in I \times I$, $j < i$ \quad $\partial_j V_i^\circ \subset
\partial V_i^\circ$ (sous-variété à bord) \quad $\partial_j V_i^\circ \to
V_j^\circ$ \emph{fibration}, fibres des variétés à bord (variétés
\add{sans bord} si $j$ prédécesseur de $i$)
\[
  \partial V_i^\circ = \bigcup_{j \in \Phi_i} \partial_j V_i^\circ
\]
\note{sous $j \in \Phi_i$ : « ens.\ de \uncertain{prédécesseurs} de $i$ »,
et une flèche de $\partial_j V_i^\circ$ vers $V_j^\circ$}
\[
  \partial(\partial_j V_i^\circ) = \bigl[ (\partial_j V_i^\circ) \,|\,
  \partial V_j^\circ \bigr] \cup \partial_{\cdot}(\partial_j V_i^\circ)
\]
\note{le second terme est souligné d'une accolade : « fibré des bords » ; son indice est illisible, \ill{}}
\begin{align*}
  \partial(\partial_j V_i^\circ) &= \bigcup_{k \in \Phi_j} (\partial_j
  V_i^\circ \,|\, \partial_k V_j^\circ) \cup \partial_{i,j}(\partial_j
  V_i^\circ) \\
  &= \Bigl( \bigcup_{k < j} (\partial_j V_i^\circ \,|\, \partial_k
  V_j^\circ) \Bigr) \cup \partial_{\emptyset,j} V_i^\circ \\
  \partial(\partial_{j'} V_i^\circ) &= \Bigl( \bigcup_{k < j'}
  \partial_{j'} V_i^\circ \,|\, \partial_k V_j^\circ \Bigr) \cup
  \partial_{i,j'} V_i^\circ
\end{align*}
\note{sous les termes, des accolades : « fibré sur $V_k^\circ$ »,
« fibré sur $V_j^\circ$ », « fibré sur $V_k^\circ$ », et pour le dernier
terme « \uncertain{vide} si $j$ prédécesseur de $i$ » ; les indices
$\partial_{i,j}$, $\partial_{\emptyset,j}$ sont d'une lecture incertaine}
\begin{align*}
  \partial(\partial_0 V_2^\circ) &\simeq \partial(\partial_1 V_2^\circ) =
  \partial_1 V_2^\circ \,|\, \partial V_1^\circ \\
  \partial_j V_i^\circ \,|\, \partial_k V_j^\circ &= \partial_{j,k}
  V_i^\circ \quad \text{fibré sur } V_k^\circ \\
  \partial_{j > k > l} &= \partial_{j,k} V_i^\circ \,|\, \partial_l
  V_k^\circ \quad \text{fibré sur } V_l^\circ \\
  & \partial_{j > k > l > m} V_i^\circ
\end{align*}

\note{Croquis : à gauche, les trois $V_0^\circ$, $V_1^\circ$, $V_2^\circ$,
avec $\partial V_1^\circ \to V_0^\circ$ et $\partial V_2^\circ =
\partial_1 V_2^\circ \cup \partial_0 V_2^\circ$ envoyé sur $V_1^\circ$ et
$V_0^\circ$ par des flèches courbes ; en bas, un tableau de flèches
$\partial_2 \leftarrow \partial_{21}$, $\partial_{20} \leftarrow
\partial_{210}$, $\partial_1 \leftarrow \partial_{10}$, $\partial_0
\leftarrow \ldots$, avec « bord » ; deux bandes hachurées étiquetées
$\partial_{210}$, $\partial_{20}$, $\partial_{21}$ ; un cylindre et un
tore étiquetés $\partial_{10}$, $\partial_1$ ; un grand cercle $\partial_0$
contenant une haltère $\partial_0$ ; dans un cadre, $\partial_1$,
$\partial_{10}$, $\partial_0 = \partial_0$, et un demi-cercle posé sur un
segment, marqué $\partial_1\ \partial_0\ \partial_1$}

\end{document}
